% !TeX program = xelatex
% !TeX encoding = UTF-8 Unicode
% !TeX TXS-program:compile = txs:///xelatex
% !TeX TXS-program:quick = txs:///xelatex | txs:///view
\documentclass[12pt,a4paper]{ctexart}
\usepackage[margin=2.25cm]{geometry}
\usepackage{amsmath}
\usepackage{booktabs}
\usepackage{tabularx}
\usepackage{array}
\usepackage{enumitem}
\usepackage{xcolor}
\usepackage{hyperref}
\usepackage{xurl}
\usepackage{listings}

\hypersetup{hidelinks}
\setlength{\parindent}{2em}
\setlength{\parskip}{0.35em}
\renewcommand{\arraystretch}{1.18}
\setlist[itemize]{leftmargin=2.2em}
\setlist[enumerate]{leftmargin=2.4em}
\lstset{
  basicstyle=\ttfamily\small,
  breaklines=true,
  columns=fullflexible,
  frame=single,
  backgroundcolor=\color{black!3}
}

\title{\textbf{面向 Ansys Maxwell 的工业软件智能体构建与验证}\\
\large 一种由自然语言驱动的仿真规格生成、脚本执行与反馈修正方法}
\author{工业软件智能体化项目阶段报告}
\date{2026年4月29日}

\begin{document}
\maketitle

\begin{abstract}
工业仿真软件通常功能强大，但使用门槛较高。用户需要理解建模、材料、激励、边界、求解器和后处理等多个环节，才能得到可信结果。本项目以 Ansys Maxwell 为对象，构建了一个工业软件智能体原型：用户输入自然语言需求后，系统将需求整理为仿真规格，生成可执行的 PyAEDT 脚本，在本地调用 Maxwell 完成建模、求解和结果提取，并根据结果判断需求是否满足。若关键约束未满足，系统会把仿真结果和失败原因反馈给大模型，由大模型生成下一轮修改方案，再重新仿真验证。

与简单的问答系统不同，本系统的输出不是文字建议，而是实际生成的 \texttt{.aedt} 工程文件、结构化结果和约束判定。当前已验证的任务族包括二维电磁铁、二维平行板电容器、二维变压器概念模型和二维电感器概念模型。实验表明，该方法能够把“中文需求”转化为“可运行的 Maxwell 仿真任务”，并在电压、电流等数值约束冲突时完成反馈修正。本文以通俗方式说明系统方法、实现流程、实验结果、复现步骤和现阶段边界。
\end{abstract}

\noindent\textbf{关键词：}工业软件智能体；Ansys Maxwell；PyAEDT；大模型；仿真自动化；反馈修正

\newpage
\tableofcontents
\newpage

\section{研究背景与问题}

工业仿真软件的核心价值在于求解真实工程问题，但它的操作链条很长。以 Maxwell 为例，一个看似简单的电磁问题，也通常需要完成几何建模、材料设置、激励设置、边界条件、求解器配置、后处理和结果判断。对不熟悉软件的人来说，真正困难的不是“点哪个按钮”，而是如何把工程需求翻译成软件能够执行的仿真任务。

本项目要解决的问题可以概括为一句话：

\begin{quote}
能否让用户用自然语言描述工程目标，由智能体自动生成并执行 Maxwell 仿真，并告诉用户结果是否满足需求？
\end{quote}

这里的“自动”不是让大模型凭空写出 Maxwell 工程文件。更可行的路线是：让大模型负责理解需求、形成仿真规格和修改策略；让本地程序负责校验、执行、调用 Maxwell 和保存工程。换句话说，大模型负责“想清楚要做什么”，Maxwell 负责“真实地算出来”，本地执行器负责“把二者可靠地连接起来”。

\section{总体思路}

\subsection{为什么不直接让大模型生成 \texttt{.aedt} 文件}

\texttt{.aedt} 是 Maxwell 的工程文件，不适合作为大模型直接生成的目标。原因有三点：

\begin{itemize}
  \item 文件结构与软件版本、内部对象和求解器状态密切相关，直接生成字节流不稳定。
  \item 即使生成了文件，也很难保证 Maxwell 能正确打开、求解和后处理。
  \item 工程文件应该由 Maxwell 自己保存，这样版本兼容性和对象关系更可靠。
\end{itemize}

因此，本项目采用间接生成方式：

\begin{center}
自然语言需求 $\rightarrow$ 仿真规格 $\rightarrow$ PyAEDT 脚本 $\rightarrow$ Maxwell 执行 $\rightarrow$ \texttt{.aedt}
\end{center}

这种路线的关键是：大模型不直接伪造工程文件，而是生成能够驱动 Maxwell 建模和求解的中间结果。

\subsection{核心闭环}

系统的主流程如下：

\begin{enumerate}
  \item 用户输入中文工程需求。
  \item 系统提取任务目标、几何参数、材料、激励、约束和需要输出的结果。
  \item 大模型生成仿真规格和执行计划。
  \item 系统生成或修复 PyAEDT 脚本。
  \item 本地调用 Maxwell 完成建模、求解和保存工程。
  \item 系统读取结果，判断每条数值约束是否满足。
  \item 若不满足，将结果和失败原因反馈给大模型，生成下一轮修改方案。
\end{enumerate}

用更形式化的方式表示，用户需求记为 $R$，系统先得到仿真规格：

\[
\Sigma = \{G, M, E, B, S, Q, C\}
\]

其中 $G$ 表示几何，$M$ 表示材料，$E$ 表示激励，$B$ 表示边界，$S$ 表示求解器，$Q$ 表示待提取结果，$C$ 表示约束。随后生成执行计划：

\[
P = \{a_1, a_2, \ldots, a_n\}
\]

其中每个 $a_i$ 是一项 Maxwell 操作，例如创建几何、赋材料、施加电流、建立求解设置或导出结果。执行后得到结果 $Y$ 和评价 $V$。如果 $V$ 表示失败，则进入反馈更新：

\[
\Sigma_{k+1}, P_{k+1} = \mathrm{LLM}(R, \Sigma_k, P_k, Y_k, V_k)
\]

当评价通过、达到最大迭代次数或无法修正时停止。

\section{系统方法}

\subsection{自然语言到仿真规格}

用户输入通常是非结构化的，例如：

\begin{quote}
做一个24V直流电磁铁，气隙2mm，电流不超过2A，尽量提高吸力，外形不要太大。
\end{quote}

这句话包含了不同层次的信息：

\begin{itemize}
  \item 数值条件：24V、2mm、不超过2A。
  \item 设计目标：尽量提高吸力。
  \item 软约束：外形不要太大。
  \item 隐含信息：需要选择电磁铁结构、线圈参数、铁芯尺寸和求解方式。
\end{itemize}

系统会把这些内容整理为统一的仿真规格。仿真规格不是最终工程文件，而是一个机器可读的任务说明。它回答的问题是：需要做什么几何、用什么材料、加什么激励、跑什么求解器、提取什么结果、检查什么约束。

\subsection{仿真规格到 PyAEDT 脚本}

仿真规格形成后，系统需要把它变成 Maxwell 能执行的动作。本项目使用 PyAEDT 作为执行接口。PyAEDT 可以用 Python 控制 AEDT/Maxwell，例如创建二维模型、设置材料、施加激励、运行求解器和保存工程。

脚本生成过程包含两层保护：

\begin{itemize}
  \item 静态检查：检查脚本是否有入口函数、是否包含 Maxwell 调用、是否存在危险操作。
  \item 运行隔离：生成脚本在独立 Python 进程中运行，避免脚本异常影响主程序。
\end{itemize}

这样做的目的不是追求形式复杂，而是保证“自动生成脚本”这件事可控。大模型可以参与生成和修复脚本，但本地系统必须负责安全边界和执行稳定性。

\subsection{结果评价}

工业仿真的重点不是“软件是否跑完”，而是“需求是否满足”。因此系统把结果分成两层判断：

\begin{itemize}
  \item 执行状态：Maxwell 是否完成建模、求解和导出。
  \item 约束状态：用户提出的数值条件是否满足。
\end{itemize}

例如在电磁铁问题中，如果用户要求 24V 且电流不超过 2A，系统不会只看 Maxwell 是否运行成功，还会估算线圈电阻和 24V 下的工作电流。若电阻为 $R_{\mathrm{coil}}$，供电电压为 $V$，则估算电流为：

\[
I_{\mathrm{est}} = \frac{V}{R_{\mathrm{coil}}}
\]

若要求 $I_{\mathrm{est}} \leq I_{\max}$，则等价于：

\[
R_{\mathrm{coil}} \geq \frac{V}{I_{\max}}
\]

这使得“24V”和“不超过2A”不再只是文本条件，而是可计算、可判定的工程约束。

\subsection{反馈修正}

当结果不满足约束时，系统不会简单报错结束，而是把失败原因反馈给大模型。例如首轮设计可能得到：

\begin{itemize}
  \item 供电电压：24V。
  \item 电流上限：2A。
  \item 估算电流：14.9A。
  \item 失败原因：24V 下电流远超上限。
\end{itemize}

大模型收到的不是一句“失败了”，而是结构化反馈，包括当前设计、仿真输出、约束判定和未满足项。随后它生成下一轮设计修改，例如增加匝数、调整线圈尺寸或改变电流设定。系统再用新设计重新生成脚本、重新运行 Maxwell、重新评价结果。

这就是当前系统区别于普通脚本自动化的核心：它不是只跑一次，而是具有“结果反馈到需求修正”的闭环。

\section{已实现任务族}

当前系统已经验证了四类任务族。这里的“验证”是指可以生成 \texttt{.aedt} 工程、执行 Maxwell、导出结果并完成评价。

\begin{table}[h]
\centering
\begin{tabularx}{\textwidth}{>{\raggedright\arraybackslash}p{3.1cm}X>{\raggedright\arraybackslash}p{3.1cm}}
\toprule
任务族 & 当前能力 & 说明 \\
\midrule
二维电磁铁 & 建模、磁静态求解、最大磁密提取、电压电流约束评价、反馈修正 & 当前最完整的闭环任务 \\
二维电容器 & 平行板几何、静电求解、电场和电容提取 & 用于验证系统跨物理场能力 \\
二维变压器 & 概念几何、匝比、电压估算、磁密提取 & 属于概念验证模型，不等于完整暂态耦合模型 \\
二维电感器 & 磁静态求解、磁密提取、等效电感估算 & 属于 Maxwell 结果与解析估算结合 \\
\bottomrule
\end{tabularx}
\end{table}

这四类任务说明系统已经从单一电磁铁任务扩展到多任务族，但还不能宣称支持任意 Maxwell 问题。更准确的表述是：系统骨架具备通用性，稳定落地能力正在逐类扩展。

\section{实验设计与结果}

\subsection{实验设置}

实验环境如下：

\begin{itemize}
  \item 操作系统：Windows。
  \item 仿真软件：Ansys Electronics Desktop Student 2025 R2。
  \item 执行接口：PyAEDT。
  \item 大模型接口：Codexa 兼容接口。
  \item 默认模型：gpt-5.4。
  \item 推理强度：high。
  \item 最大反馈迭代次数：2。
\end{itemize}

每次实验都会在 \texttt{workspace} 下生成一个独立目录。该目录保存用户需求、仿真规格、执行计划、生成脚本、脚本检查结果、Maxwell 工程文件、输出结果和评价结果。因此实验不是只依赖网页显示，而是可以从文件层面复核。

\subsection{实验一：24V 直流电磁铁}

输入需求为：

\begin{quote}
做一个24V直流电磁铁，气隙2mm，电流不超过2A，尽量提高吸力，外形不要太大。
\end{quote}

该任务的难点在于 24V 和 2A 是联合约束。系统必须保证线圈在 24V 供电时估算电流不超过 2A。最新验证结果如下：

\begin{itemize}
  \item 最大磁密：$9.866\times 10^{-5}$ T。
  \item 估算线圈电阻：13.41 $\Omega$。
  \item 24V 下估算电流：1.79 A。
  \item 总体评价：通过。
\end{itemize}

这说明系统不仅完成了 Maxwell 求解，还完成了电气约束验证。此前首轮设计曾出现电流过大的情况，经过反馈修正后，系统将设计调整到约束范围内。

\subsection{实验二：二维平行板电容器}

输入需求为：

\begin{quote}
做一个二维平行板电容器，板间距1mm，板宽20mm，施加100V电压，求电场强度和电容。
\end{quote}

结果如下：

\begin{itemize}
  \item 施加电压：100 V。
  \item 板间距：1 mm。
  \item 最大电场结果：约 $3.70\times 10^{5}$ V/m。
  \item 电容结果：198.18 pF。
  \item 总体评价：通过。
\end{itemize}

该实验的意义在于证明系统不是只围绕电磁铁场景工作。它可以进入 Maxwell 的静电求解链路，并提取电容矩阵结果。

\subsection{实验三：二维变压器概念模型}

输入需求为：

\begin{quote}
做一个10000V到220V市电的变压器。
\end{quote}

结果如下：

\begin{itemize}
  \item 原边电压：10000 V。
  \item 目标副边电压：220 V。
  \item 原边匝数：500。
  \item 副边匝数：11。
  \item 匝比：0.022。
  \item 估算副边电压：220 V。
  \item 最大磁密：$2.224\times 10^{-5}$ T。
  \item 总体评价：通过。
\end{itemize}

需要说明的是，该实验是二维概念模型，主要验证“自然语言到变压器类 Maxwell 任务”的可执行链路。它还不是完整的工频变压器高保真设计工具，暂未覆盖损耗、温升、绝缘、复杂绕组和暂态耦合等工程细节。

\subsection{实验四：二维电感器}

电感器任务用于验证系统对磁路参数和等效电感估算的处理能力。最新结果如下：

\begin{itemize}
  \item 电流：2.0 A。
  \item 线圈匝数：600。
  \item 估算电感：0.326 H。
  \item 最大磁密：$9.029\times 10^{-5}$ T。
  \item 总体评价：通过。
\end{itemize}

该任务同样属于 Maxwell 计算与简化工程估算结合的验证。它说明系统可以围绕不同输出目标组织仿真过程，而不是固定只输出磁密。

\subsection{实验结果汇总}

\begin{table}[h]
\centering
\begin{tabularx}{\textwidth}{>{\raggedright\arraybackslash}p{3.0cm}X>{\raggedright\arraybackslash}p{2.3cm}}
\toprule
实验任务 & 关键结果 & 评价 \\
\midrule
24V 直流电磁铁 & $B_{\max}=9.866\times10^{-5}$ T；$R=13.41\,\Omega$；24V 下 $I=1.79$ A & 通过 \\
二维电容器 & 电容 198.18 pF；最大电场约 $3.70\times10^{5}$ V/m & 通过 \\
二维变压器 & 匝比 0.022；估算副边电压 220 V；$B_{\max}=2.224\times10^{-5}$ T & 通过 \\
二维电感器 & 估算电感 0.326 H；$B_{\max}=9.029\times10^{-5}$ T & 通过 \\
\bottomrule
\end{tabularx}
\end{table}

\section{如何复现实验}

为了让他人能够复现，本节给出最小操作流程。这里不展开源代码，只说明运行步骤和成功判据。

\subsection{准备配置}

进入项目根目录，复制示例配置：

\begin{lstlisting}[language=bash]
cd F:\maxwell_agent_project
Copy-Item .env.example .env
\end{lstlisting}

然后在 \texttt{.env} 中填入自己的 API key，并确认模型设置为：

\begin{lstlisting}[language={}]
CODEXA_MODEL=gpt-5.4
CODEXA_REASONING_EFFORT=high
DESIGN_FEEDBACK_MAX_ITERS=2
\end{lstlisting}

\subsection{检查环境}

\begin{lstlisting}[language=bash]
.\.venv\Scripts\python.exe -m maxwell_agent.cli probe-env
.\.venv\Scripts\python.exe -m maxwell_agent.cli smoke-llm
\end{lstlisting}

第一条命令用于确认 Maxwell 和 PyAEDT 是否可用，第二条命令用于确认云端模型是否可用。

\subsection{运行电磁铁实验}

\begin{lstlisting}[language=bash]
.\.venv\Scripts\python.exe -m maxwell_agent.cli run "做一个24V直流电磁铁，气隙2mm，电流不超过2A，尽量提高吸力，外形不要太大。"
\end{lstlisting}

成功后，在新生成的 \texttt{workspace} 目录中应看到：

\begin{itemize}
  \item \texttt{electromagnet\_2d\_attempt\_*.aedt}
  \item \texttt{outputs.json}
  \item \texttt{evaluation.json}
  \item 如发生反馈修正，还会有 \texttt{intake\_feedback\_01.json} 和第二轮工程文件。
\end{itemize}

判断成功不是看命令有没有结束，而是看 \texttt{evaluation.json} 中 \texttt{overall\_status} 是否为 \texttt{passed}。

\subsection{运行电容器实验}

\begin{lstlisting}[language=bash]
.\.venv\Scripts\python.exe -m maxwell_agent.cli run "做一个二维平行板电容器，板间距1mm，板宽20mm，施加100V电压，求电场强度和电容。"
\end{lstlisting}

成功判据为：

\begin{itemize}
  \item 生成 \texttt{capacitor\_2d\_attempt\_01.aedt}。
  \item \texttt{outputs.json} 中包含电容和电场结果。
  \item \texttt{evaluation.json} 为通过。
\end{itemize}

\section{现阶段局限}

当前系统已经具备从自然语言到 Maxwell 工程文件的闭环，但仍处于原型阶段。主要局限包括：

\begin{itemize}
  \item 稳定支持的任务族仍然有限，目前主要是四类二维或概念模型。
  \item 变压器和电感器当前偏概念验证，尚未覆盖完整工程级建模细节。
  \item 复杂三维结构、旋转电机、多物理场耦合和瞬态问题仍需继续扩展执行能力。
  \item 大模型反馈修正可以提升自动化程度，但不能保证任意约束组合都一定收敛。
  \item Student 版 Maxwell 对并发和进程释放较敏感，因此当前网页端采用单任务运行策略。
\end{itemize}

这些局限并不否定当前结果。更准确地说，当前成果证明了“自然语言需求到工业仿真执行结果”的技术路线可行，但还不能把它包装成任意 Maxwell 问题都能解决的通用产品。

\section{结论}

本项目完成了一个面向 Ansys Maxwell 的工业软件智能体原型。系统能够接收中文需求，将其转化为仿真规格和执行计划，生成并检查 PyAEDT 脚本，在本地调用 Maxwell 生成 \texttt{.aedt} 工程文件，并根据结果判断约束是否满足。对于电压和电流等数值约束冲突，系统已经验证了“先仿真、再反馈、再修正、再仿真”的闭环。

从实验结果看，系统已经覆盖二维电磁铁、二维电容器、二维变压器概念模型和二维电感器概念模型四类任务。当前最重要的价值不是某一个单独算例，而是建立了一条可扩展的技术路线：把大模型用于需求理解、规格生成和反馈修正，把 Maxwell 用于真实物理求解，把本地执行器用于安全执行和结果校核。

下一步工作应集中在三方面：第一，扩展更多 Maxwell 操作原语和后处理能力；第二，提高复杂任务的反馈修正稳定性；第三，将概念模型逐步升级为更接近工程实际的高保真模型。

\end{document}
